\documentclass{ximera}
\usepackage{caption}
\graphicspath{ {./images/} }

\begin{document}
\captionsetup{singlelinecheck=false}
Directions: You have 90 minutes to complete this exam. There are 25 questions. Each question is worth 4 points. (There is no partial credit on multiple choice questions). Partial credit is rare, but possible; so please show your work. Write your answers clearly, and use the space provided for the answer!

For problems 1-5, write the letter of the correct answer in the space provided. Only one answer is correct.\\
\_1. Graph by using the slope and $y$-intercept: $y=-\frac{2}{3} x+4$

\begin{figure}[h]
\begin{center}
\captionsetup{labelformat=empty}
\caption{A)}
  \includegraphics[alt={Graph of a line with negative slope passing through (0,4) and (3,2) on a coordinate plane},width=\textwidth]{graph-choice-A.jpg}
\end{center}
\end{figure}

\begin{figure}[h]
\begin{center}
\captionsetup{labelformat=empty}
\caption{D)}
  \includegraphics[alt={Graph of a line with negative slope passing through (0,4) and (2,0) on a coordinate plane},width=\textwidth]{graph-choice-D.jpg}
\end{center}
\end{figure}

\begin{figure}[h]
\begin{center}
\captionsetup{labelformat=empty}
\caption{B)}
  \includegraphics[alt={Graph of a line with positive slope passing through (0,-4) and (3,-2) on a coordinate plane},width=\textwidth]{graph-choice-B.jpg}
\end{center}
\end{figure}

\begin{figure}[h]
\begin{center}
  \includegraphics[alt={Graph of a line with positive slope passing through (0,-4) and (2,-1) on a coordinate plane},width=\textwidth]{graph-choice-E.jpg}
\captionsetup{labelformat=empty}
\caption{E)}
\end{center}
\end{figure}

\begin{figure}[h]
\begin{center}
  \includegraphics[alt={Graph of a line with negative slope passing through (0,2) and (4,-1) on a coordinate plane},width=\textwidth]{graph-choice-C.jpg}
\captionsetup{labelformat=empty}
\caption{C)}
\end{center}
\end{figure}

$\_\_\_\_$ 2. Write the number $5.1 \times 10^{-3}$ in decimal notation.\\
A) 0.51\\
B) 0.0051\\
C) 5.1\\
D) 0.00051\\
E) 0.051\\
$\_\_\_\_$ 3. Multiply: $(4 z+5)(4 z-5)$.\\
A) $16 z^{2}+20 z+25$\\
B) $16 z^{2}-10$\\
C) $16 z^{2}+20 z-25$\\
D) $16 z^{2}-25$\\
E) $16 z^{2}+25$\\
\_ 4. Simplify: $\left(7 n^{-4} z^{-5}\right)^{2}$.\\
A) $49 n^{8} z^{10}$\\
B) $\frac{49}{n^{8} z^{5}}$\\
C) $\frac{7}{n^{8} z^{10}}$\\
D) $\frac{49}{n^{8} z^{10}}$\\
E) $\frac{49 n^{8}}{z^{10}}$\\
$\_\_\_\_$ 5. Multiply.\\
$(8 v w)(-8 v w)\left(-8 v^{3} w^{2}\right)$\\
A) $512 v^{5} w^{4}$\\
B) $-8 v^{5} w^{4}$\\
C) $-512 v^{3} w^{2}$\\
D) $-512 v^{5} w^{4}$\\
E) $512 v^{3} w^{2}$\\
6. State whether the polynomial is a monomial, binomial, or a trinomial. $4 q^{4} r^{3}-7 q^{3} r^{2}$\\
7. Simplify.\\
$\frac{5 q^{-9} r^{-2}}{-10 q r^{7}}$\\
8. Divide: $\frac{3 c^{13}-4 c^{10}-2 c^{7}}{c^{4}}$.\\
9. Find the slope of the line containing the given points.

$$P_{1}(3,-1), \quad P_{2}(8,0)$$

\begin{enumerate}
  \setcounter{enumi}{9}
  \item Subtract.
\end{enumerate}

$$\left(6 m^{3}+m-4\right)-\left(-2 m^{2}+5 m+6\right)$$

\begin{enumerate}
  \setcounter{enumi}{10}
  \item Multiply: $(v-6)^{2}$.
  \item Find the equation of the line that passes through the $y$-intercept $(0,8)$ and the point $(-5,6)$.
  \item Graph the solution set.
\end{enumerate}

$$2 x-3 y \geq-6$$

\includegraphics[alt={Blank coordinate plane with x and y axes ranging from negative 6 to 6},width=\textwidth]{blank-coordinate-plane}\\
14. Find the equation of the line that passes through the points $(-3,9)$ and $(3,-3)$.\\
15. Divide and write the answer in scientific notation.

$$\frac{1.17 \times 10^{-1}}{7.8 \times 10^{8}}$$

\begin{enumerate}
  \setcounter{enumi}{15}
  \item Solve by graphing.
\end{enumerate}

$$\begin{aligned}
3 x+5 y & =25 \\
y & =-\frac{3}{5} x+5
\end{aligned}$$

\begin{enumerate}
  \setcounter{enumi}{16}
  \item Simplify: $\left(-x^{4} y^{2}\right)^{8}$.
  \item Multiply: $\left(5 y^{3}+4 y^{2}+6\right)(6 y-1)$.
  \item Solve by substitution:
\end{enumerate}

$$\begin{aligned}
& y=2 x+13 \\
& y=3 x+18
\end{aligned}$$

\begin{enumerate}
  \setcounter{enumi}{19}
  \item Solve by the addition method.
\end{enumerate}

$$\begin{array}{r}
3 x+y=-1 \\
x-y=-1
\end{array}$$

\begin{enumerate}
  \setcounter{enumi}{20}
  \item Solve by the addition method:
\end{enumerate}

$$\begin{aligned}
& 18 x+48 y=23 \\
& 18 x+32 y=18
\end{aligned}$$

\begin{enumerate}
  \setcounter{enumi}{21}
  \item Sports A basketball team scored 70 points in two-point baskets and three-point baskets. If the two-point baskets had been three-point baskets and the three-point baskets had been two-point baskets, the team would have scored 80 points. Find how many two-point baskets and how many three-point baskets the team scored.
  \item A small plane, flying into a headwind, flew 500 miles in 7 h . Flying with the wind, the plane traveled 1150 miles in 4 h . Find the rate of the plane in calm air and the rate of the wind. Round your answer to nearest whole number.
  \item Multiply: $-3 x^{3}\left(5 x^{2}-3 x+1\right)$.
  \item Divide: $\left(2 t^{2}-2 t-10\right) \div(t-3)$.
\end{enumerate}

\section*{Answer Key}
\begin{enumerate}
  \item B
  \item B
\end{enumerate}

Grading: (4p) No partial credit.\\
:\\
3. D

Grading: (4p) No partial credit.\\
4. D

Grading: (4p) No partial credit.\\
5. A\\
6. binomial (4p) No partial credit.

\begin{center}
\begin{tabular}{l|l|l|}
7. & $\begin{array}{l}\text { (4p) } 2 \text { points for exactly one mistake, no points thereafter. Take } 1 \text { point off if the answer } \\ \text { is correct, but not fully simplified (e.g. single instances of each variable respectively, but to } \\ \text { a negative power). }\end{array}$ \\
\end{tabular}
\end{center}

\begin{enumerate}
  \setcounter{enumi}{7}
  \item $\left|3 c^{9}-4 c^{6}-2 c^{3}\right|(4 \mathrm{p}) 2$ points for exactly one computational mistake, no points thereafter.
\end{enumerate}

\begin{center}
\begin{tabular}{l|l|l|}
9. & $\frac{1}{5}$ & $\begin{array}{l}(4 \mathrm{p}) 2 \text { points for exactly one computational mistake, no points thereafter. Exception: } 1 \\ \text { point for inverse formula for the slope }(\mathrm{dx} / \mathrm{dy}) \text {, and correct numeric computation for this } \\ \text { wrong formula. }\end{array}$ \\
\hline
\end{tabular}
\end{center}

\begin{enumerate}
  \setcounter{enumi}{9}
  \item $\begin{aligned} & 6 m^{3}+2 m^{2}-4 m-\mid(4 p) 2 \text { points for exactly one computational mistake, no points thereafter. } \\ & 10\end{aligned}$
  \item $v^{2}-12 v+36 \quad \left\lvert\, \begin{aligned} & (4 \mathrm{p}) 2 \text { points for exactly one computation error; no points thereafter. } 1 \text { point if the middle } \\ & \end{aligned}\right.$ term is missing, but all other terms are correct.
  \item \begin{center}
\begin{tabular}{ll|l}
$y=\frac{2}{x+8}$ & $(4 \mathrm{p}) 2$ points for correct slope, 1 point for exactly one computational mistake; no points \\
\end{tabular}
\end{center} for wrong formula for slope. 2 points for correct equation based on found slope (even if slope is incorrect).

  \item 
\end{enumerate}

\includegraphics[alt={Graph of a linear inequality with the solution region shaded below the line on a coordinate plane},width=\textwidth]{linear-inequality-shaded-region.jpg}\\
14. $\begin{array}{ll}y=-2 x+3 & (4 p) 2 \text { points correct computation of slope; } 1 \text { point for exactly one computational mistake, }\end{array}$ no points thereafter; no points for incorrect slope formula. 2 points for correct point-slope formula, even with incorrect slope from first step; 1 point for exactly one computational mistake. Allow for not simplified answer (the question does NOT require simplification or a specific format: standard, slope-intercept).\\
15. $\left\lvert\, \begin{array}{ll}1.5 \times 10^{-10} & (4 \mathrm{p}) 2 \text { points for correct division; } 1 \text { point for exactly one computational mistake (either in }\end{array}\right.$ division of mantissa or the powers of 10); allow for minor precision errors. 2 points for correct scientific notation of answer; 1 point if the answer has the right mantissa, but the power of 10 is off; no points for any other alternatives.\\
16. infinitely many solutions (dependent)

\begin{center}
\begin{tabular}{l|l|l|}
17. & $x^{32} y^{16}$ & $\begin{array}{l}(4 \mathrm{p}) 2 \text { points for exactly one computational mistake, no points thereafter. Exception: only } \\ 1 \text { point off for wrong sign. }\end{array}$ \\
\hline
\end{tabular}
\end{center}

\begin{enumerate}
  \setcounter{enumi}{17}
  \item $30 y^{4}+19 y^{3}-4 y^{2}+36 y-6 \quad$ (4p) 2 points for exactly one computation error; no points thereafter.
  \item $(-5,3)$
  \item $\left(-\frac{1}{2}, \frac{1}{2}\right)$
  \item $\left(\frac{4}{9}, \frac{5}{16}\right)$
  \item 20 two-point baskets; 10 three-point baskets
  \item plane: 150 mph ; wind: 79 mph
  \item $-15 x^{5}+9 x^{4}-3 x^{3}$
\end{enumerate}

\begin{center}
\begin{tabular}{l|l|l}
25. & $2 t+4+\frac{2}{t-3}$ & $\begin{array}{l}(4 \mathrm{p}) 1 \text { point off for each computational mistake, up to four, no points thereafter. No } \\ \text { points for major mistakes (e.g. cancellation from sums). }\end{array}$ \\
\end{tabular}
\end{center}


\end{document}